\subsection{GitHub}
Til dette projekt brugte vi GitHub til samarbejde og fildeling. GitHub hjalp os med at organisere projektets filer, holde styr på ændringer og arbejde på forskellige dele af koden parallelt uden at risikere store konflikter. Desuden gjorde det fildeling nemmere og hjalp med at undgå, at det blev rodet.\\
Repositoryet var struktureret med en main-branch til stabile versioner og separate branches til hvert gruppemedlem for at holde main-branchen stabil. Brugen af GitHub og versionskontrol gjorde ikke kun projektarbejdet mere struktureret, men lærte os også vigtige principper for samarbejde og dokumentation i softwareudvikling.

\subsubsection{Tidsstyring og prioritering}
For at prioritere vores tid bedst muligt lavede vi både en dag-for-dag tidsplan og et GANTT-diagram for at sikre et godt overblik over vores tid. Planen blev opfattet som dynamisk og brugt som en guideline, da ændringer kunne ske dagligt. Det vigtigste udbytte af tidsplanen var at planlægge, hvordan vi skulle arbejde sammen. Det var vigtigt, at vi havde de samme prioriteringer for, hvad der var vigtigt for spillet, og sikrede, at funktionerne blev implementeret og testet i en logisk rækkefølge. Vi begyndte med at programmere spillets "need to haves" og tilføjede derefter eventuelle "nice to haves".

\subsection{Projektverificering}
Under udviklingen af spillet stødte vi på flere udfordringer og fejl, som krævede systematisk debugging for at sikre stabilitet og korrekt funktionalitet. En af de mest almindelige udfordringer var fejl i kommunikationen mellem lagene i vores program, som ofte skyldtes enten forkert initialisering af GPIO'er og timere eller fejl i logikken i vores kode.\\
For at identificere og rette logiske fejl anvendte vi forskellige debugging-teknikker, herunder breakpoints i STM32CubeIDE med variabelvisning og serielle udskrifter til terminalen. Dette gjorde det muligt at kontrollere og verificere værdierne af variabler og signaler i realtid.